% 第 5 章 IC Validator Live DRC
\bichapter{IC Validator Live DRC}{IC Validator Live DRC}
\label{chap:livedrc}

本章介绍在 Custom Compiler、IC Compiler II、Fusion Compiler 与 Cadence Virtuoso 工具中使用 IC Validator Live DRC 的方法。

IC Validator Live DRC 工具包括以下各节：
\begin{itemize}
  \item 关于 IC Validator Live DRC 工具
  \item 在 Custom Compiler 中使用 IC Validator Live DRC
  \item 在 Virtuoso 中使用 IC Validator Live DRC
  \item 在 IC Compiler II 与 Fusion Compiler 中使用 IC Validator Live DRC
  \item 设置预编译缓存目录
  \item 在 IC Validator 中复现 IC Validator Live DRC 运行
\end{itemize}

% ============================================================
\section{关于 IC Validator Live DRC 工具}
\subsection*{About the IC Validator Live DRC Tool}

先进工艺节点与大型设计上的长 DRC 运行时间会影响生产力。IC Validator Live DRC 工具使您能够直接在版图编辑器中运行 DRC，从而主动修复违例，最终得到更干净的版图并缩短签核 DRC 运行时间。

IC Validator Live DRC 工具使用 IC Validator 签核规则文件，因此检查较全面。但由于检查范围限于视口（viewport）与可见层，运行时间显著缩短。IC Validator Live DRC 工具支持在 Custom Compiler\texttrademark、IC Compiler\texttrademark\ II、Fusion Compiler\texttrademark\ 以及 Cadence\textregistered\ Virtuoso\textregistered\ 版图编辑器中使用。

% ============================================================
\section{在 Custom Compiler 中使用 IC Validator Live DRC}
\subsection*{IC Validator Live DRC With Custom Compiler}

Custom Compiler 中的 IC Validator Live DRC 允许设计人员对活动版图运行 IC Validator 物理验证检查。使用与 IC Validator 相同的 runset 与文件即可获得即时 DRC 结果。诸如 recipe 等高级选项允许创建规则子集，并可从下拉菜单中选择执行。为便于使用，违例直接显示在 Custom Compiler 中，无需单独窗口查看。以下各节说明如何设置环境、执行运行以及查看违例。

\subsection{设置环境}
\subsubsection*{Setting Up the Environment}

要在 Custom Compiler 中运行 IC Validator Live DRC，请按下列步骤操作：

\begin{enumerate}
  \item \textbf{设置环境}\\
  启动 Custom Compiler 前须完成部分 IC Validator 环境变量设置。要加载 IC Validator Live DRC 菜单，必须设置两个环境变量：
  \begin{itemize}
    \item 将 \cmd{ICV\_HOME\_DIR} 设为 IC Validator 安装目录
    \item 将 \cmd{PATH} 设为包含 \cmd{\$ICV\_HOME\_DIR/bin}
  \end{itemize}
  例如：
\begin{lstlisting}
setenv ICV_HOME_DIR </path/to/icv>
setenv PATH $ICV_HOME_DIR/bin:$PATH
\end{lstlisting}
  要使 IC Validator Live DRC 工具栏可见，选择 \textbf{Window $>$ Toolbars $>$ DRC}。

  \item \textbf{打开 IC Validator Live DRC Setup 窗口}\\
  要通过 Verification 菜单访问 IC Validator Live DRC，选择 \textbf{Tools $>$ Verification $>$ Live DRC}，或单击相应工具栏按钮。

  \item \textbf{所需文件与设置}\\
  在 Live DRC Setup 对话框中，选择下列选项卡之一：\textbf{Main}、\textbf{Control Variables}、\textbf{Custom Options}、\textbf{Include Paths}。
  \begin{enumerate}[label=\alph*.]
    \item \textbf{Main 选项卡}

\figplaceholder{Figure 18: Main Tab}{Main 选项卡}{fig:livedrc-cc-main}

\begin{longtable}{@{}>{\raggedright\arraybackslash}p{0.22\textwidth} p{0.70\textwidth}@{}}
\caption{Live DRC Setup -- Main 选项卡设置}\label{tab:livedrc-cc-main}\\
\toprule
\textbf{设置} & \textbf{说明} \\
\midrule
\endfirsthead
\caption[]{Live DRC Setup -- Main 选项卡设置（续）}\\
\toprule
\textbf{设置} & \textbf{说明} \\
\midrule
\endhead
\bottomrule
\endfoot
Run Dir & 选择运行 IC Validator Live DRC 的目录。输出文件将存放于此目录。 \\
Area & 执行 DRC 检查的区域：设计（design）、视口（viewport）或区域（region）。区域由矩形坐标（微米）按 \cmd{[x1, y1, x2, y2]} 顺序指定，例如 \cmd{[-15.5, -20.0, 15.5, 20.0]}。单击交互选择按钮可在设计中绘制选择区域。 \\
Layout Format & 对于 OpenAccess 格式，验证在当前 cell view 上运行。\\
 & 对于 Stream 与 OASIS 格式，工具导出 Stream 或 OASIS，并对生成的 GDSII 或 OASIS 文件运行 IC Validator Live DRC。请设置：
  \begin{itemize}
    \item \textbf{Format}：OpenAccess、Stream 或 OASIS——将用作输入的版图 cellview 导出到您指定的 File。
    \item 单击 \textbf{Config} 按钮打开 Export Stream 或 Export OASIS 对话框。对话框选项设置见 \textit{Custom Compiler Environment User Guide} 中 ``Exporting Design Data in GDSII Stream Format'' 或 ``Exporting Design Data in OASIS Format'' 两节。您指定的选项与 PVE（\cmd{pdk.xml} 配置中的 preferences）集成。
  \end{itemize} \\
Job Parameters & 选择下列项：
  \begin{itemize}
    \item \textbf{Runset}：用于 IC Validator Live DRC 的设计规则（runset）文件。要选择规则子集，单击相应按钮。若规则文件自加载后已修改，请重新打开对话框并单击相应按钮以加载修改后的文件。
    \item \textbf{Arguments}：控制 IC Validator Live DRC 行为的命令行参数。
  \end{itemize}
  无需选择 Tool 与 Job Class。\\
 & \textbf{Launch Debugger}：进程完成时启动（或更新已有）Error Viewer。\\
 & 在 Background 模式下，可在 IC Validator Live DRC 于后台检查时继续使用 Custom Compiler。使用 preference \cmd{xtICVLiveNotifyWhenComplete} 可在检查完成时获得通知对话框。 \\
\end{longtable}

要选择 runset（规则）文件中的特定规则，单击相应按钮。在 Live DRC Rule Selection 对话框中选择要用于规则检查的规则，然后单击 \textbf{OK}。

在过滤框中，可输入要选择的项或规则的完整或部分名称，支持简单字符串匹配、glob 风格模式或正则表达式；由 General Options 对话框中的 Filters 选项控制。要清除过滤框，单击文本字段中所示的 \textbf{X}。（见 \textit{Custom Compiler Environment User Guide} 的 ``Entering General Option Settings'' 一节。）

\figplaceholder{Figure 19: Rule Selection}{规则选择}{fig:livedrc-cc-rule-sel}

    \item \textbf{Control Variables 选项卡}\\
    在 Control Variables 选项卡中，查看并按需修改本次 DRC 运行的 DRC runset 控制变量当前值。控制变量由 PDK 预置。要修改预置值，用非选择（备用）鼠标键单击，并编辑输入字段。\\
    本对话框中的控制变量设置与 DRC Setup and Run 对话框中的设置类似。更多信息见 \textit{Design Checking and Physical Verification User Guide} 的 ``Control Variables Tab'' 与 ``Using the Integrated Tools for Design Rule Checking'' 两节。

\figplaceholder{Figure 20: Control Variables Tab}{Control Variables 选项卡}{fig:livedrc-cc-ctrl-vars}

    \item \textbf{Custom Options 选项卡}\\
    层映射文件将 OpenAccess（OA）层-purpose 名称映射到 GDSII 或 OASIS 的层号与数据类型。IC Validator Live DRC 在增量 DRC 中使用该文件。\\
    可在 Custom Options 选项卡中指定 Layer Map 文件。此外，可添加 Options 以解析 IC Validator Live DRC runset。选项列表见 \textit{Design Checking and Physical Verification User Guide} 的 ``Using the Integrated Tools for Design Rule Checking'' 一节。\\
    可使用与 stream out GDS 或 OASIS 文件时相同的层映射文件。这些文件常可从先前定义的 preferences 中取得，在 console 中输入：
\begin{lstlisting}
db::getPrefValue dbExportStreamLayerMapFile
\end{lstlisting}
    或
\begin{lstlisting}
db::getPrefValue dbExportOasisLayerMapFile
\end{lstlisting}

\figplaceholder{Figure 21: Custom Options Tab}{Custom Options 选项卡}{fig:livedrc-cc-custom-opts}

    \item \textbf{Include Paths 选项卡}\\
    在 Include Paths 选项卡中，可指定包含在 DRC 操作期间额外传递给 IC Validator Live DRC 的文件的文件夹。

\figplaceholder{Figure 22: Include Paths Tab}{Include Paths 选项卡}{fig:livedrc-cc-include}
  \end{enumerate}
\end{enumerate}

\subsection{查看作业进度}
\subsubsection*{Viewing Job Progress}

可在 job monitor 中查看 IC Validator Live DRC 作业，其中显示正确的 Status 与 Run Description。

\figplaceholder{Figure 23: Job Monitor}{作业监视器}{fig:livedrc-cc-job-monitor}

在选中的 IC Validator Live DRC 作业上右键单击，出现菜单，可：
\begin{itemize}
  \item 使用 \textbf{Kill} 命令终止正在运行的作业
  \item 查看输出文件（\cmd{icvl\_manager.log}）
\end{itemize}

\figplaceholder{Figure 24: Kill Command}{Kill 命令}{fig:livedrc-cc-kill}

若当前视图窗口中的 polygon 总数超过 50{,}000，将出现警告，因为超过该点后 IC Validator Live DRC 运行性能可能下降。若选择 ``Do not show this message again''，后续运行将继续使用 IC Validator Live DRC，除非达到最大 polygon 数限制（当前为 2{,}000{,}000）。达到该限制时无法运行 IC Validator Live DRC。

\subsection{运行 IC Validator Live DRC}
\subsubsection*{Running IC Validator Live DRC}

要启动 IC Validator Live DRC 运行：
\begin{enumerate}
  \item 在 Layout Editor 中安排版图，使要检查的区域位于视口内。
  \item 确保要检查的层可见。可用下列方法之一更改层可见性：
  \begin{itemize}
    \item 在 Hierarchy 工具栏上设置 Stop Level。
    \item 在 Object/Layer Panel 中显示或隐藏层。
  \end{itemize}
  \item 在 Live DRC 工具栏上单击运行按钮：DRC 进程启动。完成后，违例标记出现在画布上。违例也列于 Error Viewer 中。
\end{enumerate}

OA 标记显示在 Error Viewer 中（选择 \textbf{Windows $>$ Assistants $>$ Error Viewer}）。

可通过 preference \cmd{xtICVLiveNotifyWhenComplete} 控制 IC Validator Live DRC 在检查完成时是否显示通知对话框。

\figplaceholder{Figure 25: IC Validator Live DRC Console}{IC Validator Live DRC 控制台}{fig:livedrc-cc-console}

\subsection{查看与修复违例}
\subsubsection*{Viewing and Fixing Violations}

要查看 DRC 违例：
\begin{enumerate}
  \item 显示 Error Viewer。
  \item 单击 \textbf{ICV\_Live} 选项卡以查看 DRC 错误列表。
\end{enumerate}

\figplaceholder{Figure 26: Error Viewer Assistant}{Error Viewer Assistant}{fig:livedrc-cc-error-viewer}

\subsection{附加选项}
\subsubsection*{Additional Options}

诸如 halo 与 recipe 等附加选项可更好地控制运行。例如，halo 将受检数据扩展到视口窗口外紧邻区域；recipe 允许创建并运行可从下拉菜单选择的规则子集。

\subsubsection{Halo}
\paragraph*{Halos}

要在紧邻视口窗口外的数据上运行 IC Validator Live DRC，可为运行添加 halo。添加 halo 会增大检查窗口，如图~\ref{fig:livedrc-cc-halo} 所示。通过将受检数据扩展到包含视口窗口外紧邻的 polygon，可验证视口内的设计更改是否未与视口外紧邻的 polygon 产生新的 DRC 违例。

\figplaceholder{Figure 27: IC Validator Live DRC Configuration - Halo Size}{IC Validator Live DRC 配置 -- Halo 大小}{fig:livedrc-cc-halo}

黑色窗口为视口，红色窗口为视口加 halo。若以 halo 设置执行 IC Validator Live DRC：
\begin{itemize}
  \item 仅面积有一部分落在红色窗口内的 polygon 会纳入检查。（Polygon 1 纳入，Polygon 2 不纳入。）
  \item 仅与黑色窗口重叠的错误图形会纳入结果。（Error 1 纳入，Error 2 不纳入。）
\end{itemize}

可通过 preference \cmd{xtICVLiveHaloRange} 定义 halo 范围。可通过 preference \cmd{xtICVLiveDumpCheckingRange} 控制 IC Validator Live DRC 是否写出两个 OA 标记以指示检查范围：一个标记表示受检视口，另一个标记计入 halo 范围。

\subsubsection{Recipe}
\paragraph*{Recipes}

Recipe 是包含签核 runset 规则子集的 JSON 文本文件。Recipe 允许按设计阶段运行特定规则。例如，包含间距规则或前端层相关规则的 recipe 可在设计极早期执行；含更高级规则或后端层相关规则的 recipe 可在设计周期后期执行。可通过设置 Recipe File 路径创建新 recipe 或修改现有 recipe；该路径可包含多个 recipe，并用 Recipe Editor 定义。

要调用 Recipe Editor：
\begin{enumerate}
  \item 在 Live DRC Setup 对话框的 Main 选项卡中，选择现有 Recipe File 并单击相应按钮（如图~\ref{fig:livedrc-cc-recipe-args}），以加载 Recipe Editor。
  \item 单击相应按钮（如图~\ref{fig:livedrc-cc-recipe-args}），从 Select Recipe File 对话框选择 recipe。单击 \textbf{OK}，将所选 recipe 加载到 Arguments 字段。
\end{enumerate}

\figplaceholder{Figure 28: Recipe Editor With Arguments}{带 Arguments 的 Recipe Editor}{fig:livedrc-cc-recipe-args}

\begin{enumerate}
  \setcounter{enumi}{2}
  \item 选择 \cmd{<No-Recipe>}（如图~\ref{fig:livedrc-cc-recipe-ed}），以从 Arguments 路径移除 recipe 设置并创建新 recipe。
  \item 在 Rule Name 列表中勾选或取消勾选要纳入 recipe 的规则。单击相应按钮以保存新 recipe 或现有 recipe 的设置。单击相应按钮以删除指定 Recipe File 的现有 recipe。
\end{enumerate}

\figplaceholder{Figure 29: Recipe Editor}{Recipe Editor}{fig:livedrc-cc-recipe-ed}

% ============================================================
\section{在 Virtuoso 中使用 IC Validator Live DRC}
\subsection*{IC Validator Live DRC With Virtuoso}

Virtuoso 中的 IC Validator Live DRC 允许对活动版图运行 IC Validator 物理验证检查。使用与 IC Validator 相同的 runset 与文件即可获得即时 DRC 结果。诸如 recipe 等高级选项允许创建规则子集，并可从下拉菜单中选择执行。违例直接显示在 Virtuoso 中，无需单独窗口查看。以下各节说明如何设置环境、执行运行以及查看违例。

\subsection{设置环境}
\subsubsection*{Setting Up the Environment}

要在 Virtuoso 中设置 IC Validator Live DRC，须设置若干环境变量，以便 Live DRC 能定位 IC Validator 安装。
\begin{enumerate}
  \item \textbf{设置环境}\\
  将 \cmd{ICV\_HOME\_DIR} 环境变量设为 IC Validator 的安装位置。\\
  将 \cmd{SYNOPSYS\_SYSTYPE} 环境变量按系统适当设置；通常应设为 \cmd{linux64}。\\
  启动 Virtuoso，并在 CIW 窗口中输入：
\begin{lstlisting}
load( strcat( getShellEnvVar("ICV_HOME_DIR") \
"/etc/icvlive/virtuoso/icvl.il"))
\end{lstlisting}
  也可将该命令加入 \cmd{.cdsinit} 文件，以便在 Virtuoso 中自动加载 IC Validator Live DRC 工具栏。\\
  打开 Virtuoso 后，将显示 IC Validator Live DRC 工具栏，如图~\ref{fig:livedrc-v-toolbar}。可能需要拉伸工具栏才能看到全部按钮。

\figplaceholder{Figure 30: IC Validator Live DRC Toolbar}{IC Validator Live DRC 工具栏}{fig:livedrc-v-toolbar}

  \item \textbf{打开 IC Validator Live DRC Configuration 对话框}\\
  单击工具栏上的相应按钮（或键盘快捷键 \cmd{Shift+F12}）。\\
  Configuration 由三部分组成：
  \begin{itemize}
    \item \textbf{Run Settings}：确定用于检查的参数。
    \item \textbf{Check Filter}：确定将检查（或不检查）哪些层与规则。
    \item \textbf{Highlight Control}：确定浏览 IC Validator Live DRC 结果时的行为。
  \end{itemize}
\end{enumerate}

若当前视图窗口中的 polygon 总数超过 50{,}000，将出现警告，如图~\ref{fig:livedrc-v-poly-limit}。若选择 ``Do not show this message again''，后续运行将继续使用 IC Validator Live DRC，除非达到最大 polygon 数限制（当前为 2{,}000{,}000）。达到该限制时无法运行 IC Validator Live DRC。

\figplaceholder{Figure 31: Configuration - Polygon Count Limit}{配置 -- Polygon 数量限制}{fig:livedrc-v-poly-limit}

\subsection{启动 IC Validator Live DRC 检查}
\subsubsection*{Starting the IC Validator Live DRC Check}

单击相应按钮以启动 IC Validator Live DRC 检查（或键盘快捷键 \cmd{F12}）。检查期间出现模态对话框。单击 \textbf{Cancel} 可终止检查，如图~\ref{fig:livedrc-v-running}。

\figplaceholder{Figure 32: Running IC Validator Live DRC}{正在运行 IC Validator Live DRC}{fig:livedrc-v-running}

Virtuoso 命令窗口中有相关消息。IC Validator Live DRC 消息以 \cmd{[ICV Live]} 为前缀，向用户说明 Live DRC 正在做什么，并给出作业可能终止的原因信息。部分示例如图~\ref{fig:livedrc-v-msgs}。

\figplaceholder{Figure 33: IC Validator Live DRC Messages}{IC Validator Live DRC 消息}{fig:livedrc-v-msgs}

\subsection{在 IC Validator Live DRC 中浏览结果}
\subsubsection*{Browsing Results in IC Validator Live DRC}

使用 IC Validator Live DRC 工具栏可：
\begin{itemize}
  \item 选择违例规则（通过右侧的违例规则框）
  \item 浏览所选规则的错误（通过中间五个按钮）
  \item 缩放到视口（通过左侧的放大镜与坐标框）
\end{itemize}

\figplaceholder{Figure 34: IC Validator Live DRC Toolbar}{IC Validator Live DRC 工具栏}{fig:livedrc-v-toolbar2}

IC Validator Live DRC 还提供弹出式违例查看器，含违例规则的详细说明以及更多高亮选项。IC Validator Live DRC Violation Viewer 如图~\ref{fig:livedrc-v-browse}。单击 IC Validator Live DRC 工具栏中的相应图标即可打开该查看器。

\figplaceholder{Figure 35: Browsing Live DRC Results}{浏览 Live DRC 结果}{fig:livedrc-v-browse}

使用 Highlight Control 可：
\begin{itemize}
  \item \textbf{Keep Existing Highlights}\\
  若勾选，不清除先前高亮。
  \item \textbf{Follow Mode}
  \begin{itemize}
    \item \textbf{Zoom with Factor}：按因子缩放到高亮图形
    \item \textbf{Pan}：移到高亮图形中心而不改变缩放
    \item \textbf{Off}：不执行任何操作
  \end{itemize}
  \item \textbf{Auto color change}（使用 Virtuoso 的 \cmd{y0}--\cmd{y9} 层进行高亮）
  \begin{itemize}
    \item \textbf{Per Highlight}：每次高亮更换颜色
    \item \textbf{Per Rule}：仅当当前与先前高亮不属于同一规则时更换颜色
    \item \textbf{Fixed}：使用所选颜色进行高亮
  \end{itemize}
  \item Highlight Control 中的设置与 Violation Viewer 中的设置同步。
\end{itemize}

\figplaceholder{Figure 36: IC Validator Live DRC Highlight Control in the Configuration Window}{Configuration 窗口中的 Highlight Control}{fig:livedrc-v-highlight}

\subsection{附加选项}
\subsubsection*{Additional Options}

诸如 halo 与 recipe 等附加选项可更好地控制运行。例如，halo 将受检数据扩展到视口窗口外紧邻区域；recipe 允许创建并运行可从下拉菜单选择的规则子集。

\subsubsection{Halo}
\paragraph*{Halos}

添加 halo 会增大检查窗口，如图~\ref{fig:livedrc-v-halo}。通过将受检数据扩展到包含视口窗口外紧邻的 polygon，可验证视口内的设计更改是否未与视口外紧邻的 polygon 产生新的 DRC 违例。

\figplaceholder{Figure 37: IC Validator Live DRC Configuration - Halo Size}{IC Validator Live DRC 配置 -- Halo 大小}{fig:livedrc-v-halo}

黑色窗口为视口，红色窗口为视口加 halo。若以 halo 设置执行 IC Validator Live DRC：
\begin{itemize}
  \item 仅面积有一部分落在红色窗口内的 polygon 会纳入检查。（Polygon 1 纳入，Polygon 2 不纳入。）
  \item 仅与黑色窗口重叠的错误图形会纳入结果。（Error 1 纳入，Error 2 不纳入。）
\end{itemize}

\subsubsection{检查层}
\paragraph*{Check Layers}

Check Layers 选项根据可见层控制检查时纳入哪些规则。可用选项如图~\ref{fig:livedrc-v-check-layers}。

\figplaceholder{Figure 38: IC Validator Live DRC Configuration - Check Layers}{IC Validator Live DRC 配置 -- Check Layers}{fig:livedrc-v-check-layers}

各选项如下：
\begin{itemize}
  \item \textbf{All Layers}：纳入全部规则进行检查
  \item \textbf{Visible Layer Only}：仅纳入使用可见层的规则
  \item \textbf{Visible Layer Involved}：纳入至少涉及一个可见层的规则
\end{itemize}

例如，考虑下列三条规则，并假设仅 VIA 可见：
\begin{lstlisting}[style=pxl]
Rule1 @= {
       @ "Rule1";
       area(VIA, value = [2.0,2.0]);
}
Rule2 @= {
       @ "Rule2"; VIA and M1;
}
Rule3 @= {
       @ "Rule3"; M1 and M2;
}
\end{lstlisting}

选择 \textbf{All Layers} 时，纳入 Rule1、Rule2 与 Rule3。\\
选择 \textbf{Visible Layer Only} 时，仅纳入 Rule1。\\
选择 \textbf{Visible Layer Involved} 时，纳入 Rule1 与 Rule2。

\subsubsection{遵循 P\&R Boundary 可见性}
\paragraph*{Honor the P\&R Boundary visibility}

当 Check Layers 设为 \textbf{Visible Layer Only} 或 \textbf{Visible Layer Involved} 时，\textbf{Honor the P\&R Boundary visibility} 选项生效。

未设置该选项时，P\&R Boundary 对象始终视为可见。

设置该选项时，Live DRC 遵循 Virtuoso 编辑器中 P\&R Boundary 对象的可见性。

\figplaceholder{Figure 39: IC Validator Live DRC Configuration - Honor the P\&R Boundary visibility}{IC Validator Live DRC 配置 -- 遵循 P\&R Boundary 可见性}{fig:livedrc-v-pr-boundary}

\subsubsection{Recipe}
\paragraph*{Recipes}

Recipe 是包含签核 runset 规则子集的 JSON 文本文件。Recipe 允许按设计阶段运行特定规则。例如，包含间距规则或前端层相关规则的 recipe 可在设计极早期执行；含更高级规则或后端层相关规则的 recipe 可在设计周期后期执行。要在 IC Validator Live DRC 中创建 recipe，使用图~\ref{fig:livedrc-v-create-recipe} 所示功能。

\figplaceholder{Figure 40: Create New Recipe}{创建新 Recipe}{fig:livedrc-v-create-recipe}

\begin{itemize}
  \item \textbf{Rule Recipe}：定义 IC Validator Live DRC 运行中检查的规则。Live DRC 将 recipe 分为两类，如图~\ref{fig:livedrc-v-recipe-cat}：
  \begin{itemize}
    \item \textbf{User recipes}：位于 \cmd{\$HOME/.icvlive/recipes} 的 recipe。
    \item \textbf{Others recipes}：位于 \cmd{\$ICVLIVE\_RECIPE\_PATH} 的 recipe。
    \item 单击相应按钮可刷新下拉列表内容。
    \item 单击相应按钮可打开 Recipe Editor 并编辑当前所选 recipe。
  \end{itemize}

\figplaceholder{Figure 41: Recipe Categories: User and Others}{Recipe 类别：User 与 Others}{fig:livedrc-v-recipe-cat}

  \item File 菜单有三个新操作：\textbf{New}、\textbf{Rename} 与 \textbf{Delete}。

\figplaceholder{Figure 42: File Menu}{File 菜单}{fig:livedrc-v-file-menu}

  \item Recipes 菜单列出全部可用 recipe，并允许在 recipe 之间切换。

\figplaceholder{Figure 43: Recipes}{Recipes}{fig:livedrc-v-recipes-menu}

  \item Selected 与 Unselected 规则列表：
  \begin{itemize}
    \item 每组规则总数显示在表头括号中。
    \item 要在 Selected 与 Unselected 规则列表之间移动规则：
    \begin{itemize}
      \item 单击列表项以选择规则
      \item 支持连续选择（按 \cmd{Shift}）、扩展选择（按 \cmd{Ctrl}）与全选（\cmd{Ctrl+a}）
      \item 单击相应按钮以移动规则项
    \end{itemize}
  \end{itemize}

\figplaceholder{Figure 44: Recipe Creating Rules}{Recipe 创建规则}{fig:livedrc-v-recipe-rules}
\end{itemize}

\subsubsection{保存与加载配置}
\paragraph*{Saving and Loading Configurations}

配置文件允许保存 IC Validator Live DRC 运行所用的 runset 及其他设置。保存设置后，即使关闭先前会话并打开新版图或会话，仍可恢复工作。配置文件是与 IC Validator Recipe 文件类似的基于 JSON 的文本文件。本节说明如何创建这些文件并将其加载到 IC Validator Live DRC 运行。要在运行前加载配置文件，设置 \cmd{ICV\_LIVE\_CONFIG\_FILE}，后跟以冒号分隔的文件或目录。例如：
\begin{lstlisting}
setenv ICV_LIVE_CONFIG_FILE file1:/path/to/dir1:file2:$cwd
\end{lstlisting}

对具有写权限的任何文件，可直接从 IC Validator Live DRC Configuration Window 编辑。

设置好配置后，可单击保存按钮并使用弹出对话框进行保存。

若该配置常用，自动加载会很有用。可通过环境变量 \cmd{ICV\_LIVE\_CONFIG\_FILE} 实现：将其设为以冒号分隔的文件或目录列表。打开 IC Validator Live DRC 配置菜单时，Live DRC 读取 \cmd{ICV\_LIVE\_CONFIG\_FILE} 路径中的配置文件。随后可从 Configuration 菜单顶部的配置下拉列表中选择配置。

\figplaceholder{Figure 45: User Configuration}{用户配置}{fig:livedrc-v-user-cfg}

添加 runset 与层映射。Configuration 窗口顶部显示的配置下拉菜单可用于在不同配置设置之间切换。

\paragraph{保存配置}
\subparagraph*{Saving Configurations}

要保存配置：
\begin{enumerate}
  \item 在 ICV Live Configuration 中，单击 \textbf{OK} 或 \textbf{Apply}。
  \item 使用 Configuration 下拉列表切换到另一配置。
\end{enumerate}

\figplaceholder{Figure 46: Saving Configurations}{保存配置}{fig:livedrc-v-save-cfg}

\begin{enumerate}
  \setcounter{enumi}{2}
  \item 单击 \textbf{Create Configuration} 或 \textbf{Load Configuration} 文件。
\end{enumerate}

\figplaceholder{Figure 47: Create Configuration}{创建配置}{fig:livedrc-v-create-cfg}

IC Validator Live DRC 按如下方式检查配置的写权限：
\begin{enumerate}[label=\alph*.]
  \item 若未保存的配置可写，Live DRC 将更改保存到可写的配置文件。
  \item 若未保存的配置为只读，则出现提示，要求将配置保存到新的可写配置文件，或覆盖现有可写配置文件。
\end{enumerate}

\figplaceholder{Figure 48: User Configuration}{用户配置}{fig:livedrc-v-user-cfg2}

\figplaceholder{Figure 49: Loading an Existing Configuration File}{加载现有配置文件}{fig:livedrc-v-load-cfg}

\paragraph{配置文件内容}
\subparagraph*{Contents of a Configuration File}

配置文件以 JSON 格式编写。IC Validator Live DRC 中的配置文件用于在后续新会话中记住初始用户设置。一个配置文件可包含多个配置。例如，某 N7 配置文件可能有两个配置——一个用于 BEOL 运行，一个用于 FEOL 运行。

例如：
\begin{lstlisting}
%vi n7_configurations.icvlc
\end{lstlisting}

\begin{lstlisting}[style=shell]
{
    "configurations": [
        {
            "checkFilter": {
                "checkLayers": "Visible Layers Involved",
                "excludeDensityRules": true,
                "loadedRecipeFiles": [
                    "/path/to/beol_spacing.icvrcp",
                    "/path/to/beol_density.icvrcp"
                ],
                "selectedRecipe": "Spacing Rules",
                "selectedRecipeFile": "/path/to/beol_spacing.icvrcp"
            },
            "name": "BEOL Configuration",
            "runSettings": {
                "clfFile": "/path/to/beol_clf.txt",
                "enableColoring": true,
                "exportLayoutToOasis": true,
                "haloSize": 1.5,
                "layerMappingFile": "n7.layermap",
                "objectMappingFile": "n7.objmap",
                "runsetFile": "n7.rs"
            }
        },
        {
            "checkFilter": {
                "checkLayers": "All Layers",
                "excludeDensityRules": true,
                "loadedRecipeFiles": [
                ],
                "selectedRecipe": "",
                "selectedRecipeFile": ""
            },
            "name": "All Layers ",
            "runSettings": {
                "clfFile": "",
                "enableColoring": false,
                "exportLayoutToOasis": false,
                "haloSize": 1,
                "layerMappingFile": "n7.layermap",
                "objectMappingFile": "n7.objmap",
                "runsetFile": "n7.rs"
            }
        },
        {
            "checkFilter": {
                "checkLayers": "All Layers",
                "excludeDensityRules": true,
                "loadedRecipeFiles": [
                ],
                "selectedRecipe": "Custom Checks",
                "selectedRecipeFile": "/path/to/feol.icvrcp"
            },
            "name": "FEOL Configuration",
            "runSettings": {
                "clfFile": "/path/to/feol_clf.txt",
                "enableColoring": true,
                "exportLayoutToOasis": true,
                "haloSize": 1.5,
                "layerMappingFile": "n7.layermap",
                "objectMappingFile": "n7.objmap",
                "runsetFile": "n7.rs"
            }
        }
    ],
    "majorVersion": 1,
    "minorVersion": 0
}
\end{lstlisting}

\subsubsection{Preference 设置}
\paragraph*{Preference Settings}

Live DRC 将与 preference 相关的设置存入每用户文件。这些设置不影响 Live DRC 检查本身，但影响 GUI 与高亮。

例如，Highlight Control 部分存为 preference。

默认情况下，该信息存于 \cmd{\$HOME/.synopsys/icvlive/prefs.icvlc}。
要更改存储文件，设置 \cmd{ICV\_LIVE\_PREF\_FILE}，后跟文件名。
例如：\cmd{setenv ICV\_LIVE\_PREF\_FILE /path/to/my\_prefs.icvlc}

\paragraph{运行目录}
\subparagraph*{Run Directory}

设置环境变量 \cmd{\_ICV\_LIVE\_CLIENT\_ENABLE\_GUI\_RUN\_DIR} 可启用如图~\ref{fig:livedrc-v-run-dir} 所示的 Run Directory 字段。用该字段指定 Live DRC 写入中间文件的目录。未指定目录时，Live DRC 默认将中间文件保存到 \cmd{\$HOME/.synopsys/icvlive}。

\figplaceholder{Figure 50: Run Directory Field In Preference Settings}{Preference 设置中的 Run Directory 字段}{fig:livedrc-v-run-dir}

要覆盖 GUI 设置，设置环境变量 \cmd{ICV\_LIVE\_OUTPUT\_DIR}。该环境变量优先级高于 GUI 中的 Run Directory 字段。

\subsubsection{快捷键}
\paragraph*{Bindkeys}

IC Validator Live DRC 快捷键可快速访问常用功能。有两个默认可写在 \cmd{icvl.il} 文件末尾的快捷键，该文件位于 \cmd{\$ICV\_HOME\_DIR/etc/icvlive/virtuoso/icvl.il}。

\begin{itemize}
  \item \cmd{Shift+F12}：打开配置窗口。
  \item \cmd{F12}：运行 IC Validator Live DRC。
\end{itemize}

可通过 \cmd{hiSetBindKey()} 函数定义快捷键，以设置 Layout、Key 以及要调用的 SKILL 函数，如下所示：
\begin{itemize}
  \item 填入 \cmd{icvlSetup(hiGetCurrentWindow())} 以打开配置窗口。
  \item 填入 \cmd{icvlExecute(hiGetCurrentWindow())} 以运行 IC Validator Live DRC。
  \item 填入 \cmd{icvlBrowse(hiGetCurrentWindow())} 以打开违例查看器。
\end{itemize}

例如：
\begin{lstlisting}
# Set hotkey for IC Validator Live DRC configuration dialog and execution
hiSetBindKey("Layout" "Shift<Key>F12" "icvlSetup(hiGetCurrentWindow())")
hiSetBindKey("Layout" "<Key>F12" "icvlExecute(hiGetCurrentWindow())")
\end{lstlisting}

% ============================================================
\section{在 IC Compiler II 与 Fusion Compiler 中使用 IC Validator Live DRC}
\subsection*{IC Validator Live DRC With IC Compiler II and Fusion Compiler}

可在 IC Compiler II 或 Fusion Compiler GUI 中使用 IC Validator Live DRC 工具执行交互式签核设计规则检查。

\subsection{设置环境}
\subsubsection*{Setting Up the Environment}

要使用 IC Validator Live DRC 执行交互式签核设计规则检查：
\begin{enumerate}
  \item 设置 IC Validator 环境。通过设置 \cmd{ICV\_HOME\_DIR} 环境变量指定 IC Validator 可执行文件的位置。可在 \cmd{.cshrc} 文件中设置该变量。使用类似下列示例的命令指定位置：
\begin{lstlisting}
setenv ICV_HOME_DIR /root_dir/icv
set path = ($path $ICV_HOME_DIR/bin)
\end{lstlisting}
  须确保指定的 IC Validator 可执行文件版本与所用的 IC Compiler II 或 Fusion Compiler 版本兼容。使用 IC Compiler II 或 Fusion Compiler 的 \cmd{report\_versions} 命令可报告版本兼容性。

  \item 设置交互式签核设计规则检查的应用选项。\\
  至少须通过设置 \cmd{signoff.check\_drc.runset} 应用选项指定用于设计规则检查的 foundry runset，并通过设置 \cmd{signoff.physical.layer\_map\_file} 应用选项指定将技术文件层映射到 runset 层的层映射文件。\\
  运行交互式设计规则检查前，按表~\ref{tab:livedrc-icc2-appopts} 设置交互式签核设计规则检查的应用选项以配置运行。使用 \cmd{set\_app\_options} 命令设置应用选项；使用 \cmd{report\_app\_options} 命令查看当前设置。
\end{enumerate}

\begin{longtable}{@{}>{\raggedright\arraybackslash\ttfamily}p{0.36\textwidth} >{\raggedright\arraybackslash}p{0.16\textwidth} p{0.40\textwidth}@{}}
\caption{交互式设计规则检查的应用选项}\label{tab:livedrc-icc2-appopts}\\
\toprule
\textbf{\textrm{应用选项}} & \textbf{\textrm{默认值}} & \textbf{\textrm{说明}} \\
\midrule
\endfirsthead
\caption[]{交互式设计规则检查的应用选项（续）}\\
\toprule
\textbf{\textrm{应用选项}} & \textbf{\textrm{默认值}} & \textbf{\textrm{说明}} \\
\midrule
\endhead
\bottomrule
\endfoot
signoff.check\_drc\_live.default\_layers & (none) & 指定写入 IC Validator Live DRC 默认层的带 color mask 的层。 \\
signoff.check\_drc\_live.default\_halo & 1 & 指定以微米为单位扩展 IC Validator Live DRC 执行检查之区域的距离。 \\
signoff.check\_drc\_live.keep\_enclosing\_results & false & 指定是否保留检查区域外的 DRC 违例。默认情况下，IC Validator Live DRC 仅保留完全位于检查区域内的违例。若将该应用选项设为 \cmd{true}，则保留完全位于 trim 区域内的 DRC 违例。trim 区域为检查区域加上 \cmd{signoff.check\_drc\_live.halo} 应用选项所指定距离的一半。完全位于 trim 区域外的违例始终丢弃。 \\
signoff.check\_drc\_live.run\_dir & signoff\_check\_drc\_live\_run & 指定运行目录，其中包含交互式签核设计规则检查生成的文件。可指定相对路径（此时目录在当前工作目录下创建）或绝对路径。 \\
signoff.check\_drc\_live.runset\newline(required) & (none) & 指定用于设计规则检查的 foundry runset。 \\
signoff.check\_drc\_live.user\_defined\_options & (none) & 指定 IC Validator 命令行的附加选项。该选项中指定的字符串会追加到用于调用 IC Validator 工具的命令行。IC Compiler II/Fusion Compiler 工具不对指定字符串执行任何检查。 \\
signoff.physical.layer\_map\_file\newline(required) & (none) & 指定将技术文件中的层名映射到 foundry runset 中层名的层映射文件名。 \\
\end{longtable}

\begin{itemize}
  \item 确保版图窗口显示要对之执行设计规则检查的对象。\\
  默认情况下，交互式设计规则检查对每个 cell 使用 frame view。若需要比 frame view 更详细的内容，可对特定 cell 使用 design 或 layout view。使用 \cmd{gui\_set\_display\_view} 命令指定特定 cell 所用的 view。例如，要对所有层次 cell 使用 layout view：
\begin{lstlisting}
fc2_shell> gui_set_display_view \
    -cells [get_cells -hierarchical *] -view layout
\end{lstlisting}
  关于控制用于设计规则检查的显示对象的信息，见「为设计规则检查显示对象」。

  \item 在 GUI 菜单栏上右键单击并选择 \textbf{DRC Toolbar}，以显示 DRC 工具栏，如图~\ref{fig:livedrc-fc-toolbar}。

\figplaceholder{Figure 51: DRC Toolbar}{DRC 工具栏}{fig:livedrc-fc-toolbar}

  DRC Toolbar 图标见表~\ref{tab:livedrc-drc-toolbar}。

\begin{longtable}{@{}>{\raggedright\arraybackslash}p{0.48\textwidth} p{0.44\textwidth}@{}}
\caption{DRC 工具栏按钮}\label{tab:livedrc-drc-toolbar}\\
\toprule
\textbf{功能} & \textbf{说明} \\
\midrule
\endfirsthead
\caption[]{DRC 工具栏按钮（续）}\\
\toprule
\textbf{功能} & \textbf{说明} \\
\midrule
\endhead
\bottomrule
\endfoot
运行检查 & 使用 IC Validator Live DRC 运行设计规则检查。 \\
应用选项 & 打开对话框以查看并设置 IC Validator Live DRC 应用选项。 \\
配置规则 & 打开对话框以配置用于设计规则检查的规则。 \\
选择区域 & 选择执行设计规则检查的区域。 \\
重置区域 & 将设计规则检查区域重置为先前运行所用区域。 \\
清除/高亮全部违例 & 清除或高亮 IC Validator Live DRC 检测到的全部 DRC 违例。 \\
遍历违例 & 遍历当前规则来自 IC Validator Live DRC 的 DRC 违例。 \\
缩放/平移控制 & 控制遍历 DRC 违例时的缩放或平移。 \\
显示当前规则 & 显示正在高亮的设计规则。 \\
\end{longtable}

  \item 通过选择 \textbf{Edit $>$ Options $>$ ICV Live}，或单击 DRC 工具栏中的齿轮图标，配置用于设计规则检查的规则。
  \item 通过选择 \textbf{Edit $>$ ICV $>$ Run ICV-Live on Current View}，或单击 DRC 工具栏中的 ICV Run 图标，运行设计规则检查。\\
  首次运行 IC Validator Live DRC（或在更改层或规则之后），工具会缓存 runset，可能需要几分钟。后续运行不再执行该缓存，因此快得多。
\end{itemize}

\subsection{为设计规则检查显示对象}
\subsubsection*{Displaying Objects for Design Rule Checking}

运行交互式设计规则检查时，IC Validator 工具对版图窗口中显示的内容加上 1 微米扩展执行设计规则检查。

可通过设置 \cmd{signoff.check\_drc\_live.halo} 应用选项修改扩展大小。使用 View Settings 与 Hierarchy Settings 面板控制版图窗口中显示的层、层次与 cell 类型。
\begin{itemize}
  \item 单击面板区的 View Settings 图标以显示 View Settings 面板。
  \begin{itemize}
    \item 在 Level 字段中指定要显示的层级数。
    \item 选择 Layers 选项卡并为所需层启用可见性。
  \end{itemize}
  \item 单击 View Settings 面板或菜单栏中的 Hierarchy Settings 图标以显示 Hierarchy Settings 面板。
  \begin{itemize}
    \item 在 ``View level'' 字段中指定视图层级数。
    \item 在 ``Expanding cell types'' 部分选择 Standard、``Fill Cell'' 与 Block。
  \end{itemize}
\end{itemize}

图~\ref{fig:livedrc-fc-view-hier} 显示用于控制设计规则检查之版图窗口显示的 View Settings 与 Hierarchy Settings 面板。

\figplaceholder{Figure 52: View and Hierarchy Settings}{View 与 Hierarchy 设置}{fig:livedrc-fc-view-hier}

\begin{noteBox}
关于交互式签核设计规则检查的更多信息，见 SolvNetPlus 上的 \textit{Fusion Compiler\texttrademark\ User Guide} 或 \textit{IC Compiler\texttrademark\ II Implementation User Guide}。
\end{noteBox}

% ============================================================
\section{设置预编译缓存目录}
\subsection*{Setting Up Pre-Compiled Cache Directories}

每当 IC Validator Live DRC 在视图窗口中检测到新层或规则选择发生变化时，都会重新编译 runset。若频繁更改视图窗口与可见层，会导致频繁编译，从而减慢 Live DRC 运行时间。为提升性能，可准备包含设计过程中生成的缓存文件的缓存目录。当窗口日后更改时，若层、规则选择与设置与当前运行匹配，Live DRC 可从该目录搜索并重用较旧的缓存文件。这些缓存目录可与其他用户共享。要创建 runset 缓存目录并重用或共享文件：
\begin{enumerate}
  \item 设置本地缓存目录并运行 IC Validator Live DRC，以不同配置编译 runset。
\begin{lstlisting}
unset ICV_DISABLE_RUNSET_CACHE
setenv ICV_RUNSET_CACHE_DIR <first_user local>
# Start ICV Live
\end{lstlisting}
  \item 将缓存文件从 \cmd{<first\_user local>} 传输到 \cmd{<new\_user>} 的共享缓存目录 \cmd{<new\_user shared>}。
  \item 在 \cmd{<new\_user>} 中同时设置本地与共享缓存目录以运行 IC Validator Live DRC，使其既能拾取预编译文件，也能将新创建的缓存文件存入本地路径。
\begin{lstlisting}
unset ICV_DISABLE_RUNSET_CACHE
setenv ICV_RUNSET_CACHE_DIR <new_user local>
setenv ICV_RUNSET_CACHE_DIR_SHARED <new_user shared>
# Start ICV Live
\end{lstlisting}
\end{enumerate}

% ============================================================
\section{在 IC Validator 中复现 IC Validator Live DRC 运行}
\subsection*{Reproducing an IC Validator Live DRC Run in IC Validator}

要在 IC Validator 中复现 IC Validator Live DRC 运行结果：
\begin{enumerate}
  \item 从 Console 捕获运行的坐标：

\figplaceholder{Figure 53: IC Validator Live DRC Console}{IC Validator Live DRC 控制台}{fig:livedrc-reproduce-console}

  \item 在运行目录中用下列内容创建 \cmd{config.rs} 文件，并将 \cmd{select\_window} 中的参数替换为上一步 IC Validator Live DRC 运行的坐标：
\begin{lstlisting}[style=pxl]
hierarchy_options(
   explode_all = {"*"}
);
incremental_options(
   select_window = {{186.62, -33.817, 192.58, -31.087}},
   clip_window = false,
   window_error_filter = KEEP_ENCLOSED,
   window_ambit = double
);
\end{lstlisting}
  IC Validator Live DRC 在扁平数据上运行。因此必须使用 \cmd{hierarchy\_options()} 函数将输入扁平化到 IC Validator 运行。尽管 Live DRC 不在视口窗口边缘裁剪 polygon，但也不会报告视口外的任何违例。因此必须对 \cmd{incremental\_options()} 函数使用 \cmd{clip\_window} 与 \cmd{window\_error\_filter} 参数。
  \item 通过设置 \cmd{-runset\_config config.rs} 将该文件加入 IC Validator 运行。
  \item 运行 IC Validator。
\end{enumerate}
